\section{Исследование производительности и качества}

\subsection{Дневник выполнения работы}

В соответствии с заданием, ниже представлен протокол выполнения работы по дням.

\begin{center}
\begin{tabular}{|p{3cm}|p{7cm}|p{4cm}|}
\hline
\textbf{Дата} & \textbf{Описание проделанной работы} & \textbf{Результат} \\ \hline
15.05.2026 & Изучение условия задачи и выбор структуры данных для поиска образцов в заранее известном тексте. & В качестве основной структуры данных выбран суффиксный массив. \\ \hline
16.05.2026 & Реализация построения суффиксного массива методом удвоения. & Получена первая рабочая версия программы, выполняющая поиск образцов с помощью бинарного поиска по суффиксному массиву. \\ \hline
17.05.2026 & Добавление массива LCP и структуры RMQ для ускорения поиска. Реализован LCP-ускоренный бинарный поиск. & Поиск образца оптимизирован до оценки \(O(m + \log n)\), где \(m\)~--- длина образца, \(n\)~--- длина текста. \\ \hline
18.05.2026 & Проведение экспериментального исследования времени работы поиска при увеличении длины текста. & Получены результаты, подтверждающие медленный, близкий к логарифмическому, рост времени поиска при фиксированной длине образца. \\ \hline
\end{tabular}
\end{center}

\subsection{Описание проведённого бенчмарка}

Для исследования производительности была написана отдельная программа-бенчмарк. 
Она строит суффиксный массив, массив LCP и структуру RMQ, после чего выполняет серию поисковых запросов.

В данном эксперименте длина образца была зафиксирована:

\[
m = 64.
\]

Длина текста \(n\) увеличивалась в два раза на каждом шаге:

\[
10000,\ 20000,\ 40000,\ \ldots,\ 1280000.
\]

Для каждого значения \(n\) выполнялось \(10000\) запросов. 
Время вывода и сортировки позиций не учитывалось, так как исследовалась именно сложность поиска диапазона суффиксов, начинающихся с заданного образца.

Ожидаемая сложность поиска одного образца:

\[
O(m + \log n).
\]

Так как \(m\) фиксировано, при увеличении \(n\) время поиска должно расти медленно, примерно как \(\log n\).

\subsection{Результаты замеров}

Ниже приведены результаты запуска бенчмарка командой:

\begin{verbatim}
make by-n
\end{verbatim}

\begin{center}
\begin{tabular}{|r|r|r|r|r|r|r|}
\hline
\textbf{\(n\)} & 
\textbf{build, ms} & 
\textbf{search, ms} & 
\textbf{\(\mu s/query\)} & 
\textbf{\(\log_2 n\)} & 
\textbf{\(m + \log_2 n\)} & 
\textbf{\(\frac{\mu s}{m+\log n}\)} \\ \hline

10000   & 0.883   & 1.179 & 0.117873 & 13.288 & 77.288 & 0.00152512 \\ \hline
20000   & 1.804   & 1.341 & 0.134118 & 14.288 & 78.288 & 0.00171314 \\ \hline
40000   & 3.792   & 1.443 & 0.144306 & 15.288 & 79.288 & 0.00182003 \\ \hline
80000   & 8.011   & 1.440 & 0.143966 & 16.288 & 80.288 & 0.00179313 \\ \hline
160000  & 17.462  & 1.496 & 0.149647 & 17.288 & 81.288 & 0.00184096 \\ \hline
320000  & 37.063  & 1.510 & 0.150978 & 18.288 & 82.288 & 0.00183476 \\ \hline
640000  & 80.995  & 1.544 & 0.154412 & 19.288 & 83.288 & 0.00185396 \\ \hline
1280000 & 171.639 & 1.564 & 0.156424 & 20.288 & 84.288 & 0.00185583 \\ \hline
\end{tabular}
\end{center}

\subsection{Анализ результатов}

По результатам эксперимента видно, что время построения структур данных растёт значительно быстрее времени поиска. 
Это ожидаемо, так как суффиксный массив строится методом удвоения, имеющим сложность:

\[
O(n \log n).
\]

Дополнительно строятся массив LCP и структура RMQ. 
Эти этапы выполняются один раз для всего текста, после чего программа может обрабатывать множество поисковых запросов.

Наиболее важной частью эксперимента является время поиска. 
При увеличении длины текста с \(10000\) до \(1280000\), то есть в \(128\) раз, среднее время одного запроса увеличилось только с

\[
0.117873\ \mu s
\]

до

\[
0.156424\ \mu s.
\]

Таким образом, время одного поиска увеличилось примерно в

\[
\frac{0.156424}{0.117873} \approx 1.33
\]

раза. Если бы поиск зависел от длины текста линейно, то при таком увеличении размера входных данных время должно было бы вырасти примерно в \(128\) раз. 
На практике наблюдается намного более медленный рост, что соответствует логарифмической зависимости от \(n\).

Также была рассчитана величина:

\[
\frac{T}{m + \log_2 n}.
\]

Начиная с \(n = 40000\), это значение остаётся практически постоянным:

\[
0.00182003,\ 0.00179313,\ 0.00184096,\ 0.00183476,\ 0.00185396,\ 0.00185583.
\]

Это показывает, что экспериментальное время поиска хорошо согласуется с моделью:

\[
T(n, m) = C \cdot (m + \log n).
\]

\subsection{Выводы о производительности}

В ходе исследования были получены следующие результаты:

\begin{enumerate}
    \item Построение суффиксного массива, массива LCP и RMQ требует заметного времени, однако выполняется только один раз для заранее известного текста.
    
    \item После построения вспомогательных структур поиск каждого образца выполняется быстро.
    
    \item При фиксированной длине образца \(m = 64\) и увеличении длины текста \(n\) в \(128\) раз среднее время одного запроса увеличилось только примерно в \(1.33\) раза.
    
    \item Отношение \(\frac{T}{m + \log_2 n}\) остаётся почти постоянным, что экспериментально подтверждает оценку времени поиска \(O(m + \log n)\).
    
    \item Использование LCP и RMQ позволяет избежать повторного сравнения одних и тех же префиксов при бинарном поиске по суффиксному массиву.
\end{enumerate}

\subsection{Выводы о найденных недочётах}

В процессе разработки были выявлены следующие особенности и недочёты первоначального подхода:

\begin{enumerate}
    \item \textbf{Обычный бинарный поиск по суффиксному массиву}. 
    В первой версии сравнение образца с суффиксом могло каждый раз начинаться с первого символа. 
    В худшем случае это приводило к оценке \(O(m \log n)\) на один запрос.

    \item \textbf{Недостаточное использование информации о соседних суффиксах}. 
    Суффиксный массив сам по себе хранит только порядок суффиксов, но не хранит информацию об их общих префиксах. 
    Для устранения этого недостатка был добавлен массив LCP.

    \item \textbf{Необходимость быстро получать LCP произвольных суффиксов}. 
    Для LCP-ускоренного бинарного поиска требуется быстро находить наибольший общий префикс двух суффиксов по их позициям в суффиксном массиве. 
    Для этого над массивом LCP была построена структура RMQ.

    \item \textbf{Большое количество операций вывода}. 
    Так как для одного образца может существовать много вхождений, вывод позиций может быть объёмным. 
    Для уменьшения накладных расходов вывод накапливается в строковом буфере.
\end{enumerate}

\pagebreak